\section{问题分析}

\subsection{数据口径的三处关键确认}

\textbf{功率与电量的换算。}附件给出的电价单位是元/$\mathrm{kWh}$，负载与光伏单位是 $\mathrm{kW}$，而结果表要求填写 $\mathrm{kWh}$。10 分钟区间内功率近似恒定，故区间电量等于功率乘以 $1/6\,\mathrm{h}$。储能 $5000\,\mathrm{kW}$ 的功率上限对应每区间 $5000/6\approx833.333\,\mathrm{kWh}$ 的能量上限。

\textbf{附件 3 的时间对齐。}附件 3 每行是一次预报，“预报 $h$ 小时”指发布时刻起第 $h$ 个小时，即区间 $(\text{发布时刻}+h-1,\ \text{发布时刻}+h]$。以 2025-06-21 为例，0:00 发布的“预报 6 小时”为 $2533\,\mathrm{kWh}$，而当天 5:00--6:00 的实际光伏为 $1905\,\mathrm{kWh}$；“预报 6 小时”若解释成 6:00--7:00 则与实际 $3405\,\mathrm{kWh}$ 相差更大。逐小时相关分析也支持该对齐：按此对齐，4:00--17:00 各小时的预报与实际相关系数均在 $0.84$ 以上。使用时在该小时的 6 个十分钟区间内取常值。

\textbf{附件 5 模板的表头错位。}模板“计划购电量”工作表的第一行标签是 \texttt{0:10-0:20}，最后一行是 \texttt{0:00+1-0:10+1}，比当天实际的 144 个区间整体错开一格（当天的第一个区间应为 0:00--0:10）。由于题目要求不得改动模板结构，本文按行号映射：第 2 行填写第 0 个区间，第 145 行填写第 143 个区间，并在结果报告中披露该风险。附件 1、附件 2、附件 4 的“时间”列同样是区间结束时刻（0:10 至 0:00+1），与本映射自洽。

\subsection{各问的思路}

\textbf{问题一}是确定性优化：电价、负载、光伏预测全部给定，储能首尾电量相等。它本质是一个带储能递推约束的线性规划\cite{boyd2004}，可以用单纯形法一次求解。

\textbf{问题二}引入两类不确定性。其一是负载与光伏的真实值偏离预测值；其二是紧急购电的价格高达正常电价的 5 倍。题面没有给出问题二的日前预测数据，因此必须自行构造一个只使用决策时刻以前信息的因果预测\cite{hyndman2021}。取此前 7 天的历史均值作为预测，同时必须处理 $5$ 倍非对称代价：买多了只能弃掉（损失 $1$ 倍电价），买少了要按 $5$ 倍价补购，所以最优计划不应等于预测值，而应在其上叠加一个安全裕度。

\textbf{问题三}提供了更准的光伏信息。要求把附件 3 的整点预报展开到 10 分钟粒度，在四个预报时刻滚动重排尚未执行的时段，并按“增购 1.5 倍、减购 0.5 倍”结算偏差。问题的落脚点是：这些后续预报究竟值多少钱。

\textbf{问题四}把电价也变成随机量。电价同样只能在 0:00 之前观测，因此计划用历史均值价制定、按实际价结算，两者的差额就是价格预测误差的代价。

\subsection{贯穿四问的一个恒等式}

在“储能严格按计划充放电”的执行口径下，把计划的功率平衡式
$x_t+PV_t^{\text{fc}}+d_t-c_t-w_t=L_t^{\text{fc}}$
代入执行缺额 $\text{gap}_t=L_t+c_t-PV_t-d_t-x_t$，得到
\begin{equation}
\text{gap}_t=e_t-w_t,\qquad
e_t=\bigl(L_t-L_t^{\text{fc}}\bigr)-\bigl(PV_t-PV_t^{\text{fc}}\bigr),
\label{eq:identity}
\end{equation}
其中 $e_t$ 是净负荷的预测误差。\textbf{缺额与储能的充放电决策无关}，只由预测误差和计划内弃光量决定。这个恒等式有两个直接推论：第一，问题二中紧急购电只能靠“多买”（即 $w_t>0$）来压缩，不能靠储能调度；第二，问题三压缩的是同一个 $e_t$，因此预报精度的提升可以直接换算成紧急购电费用的下降。全文的问题二至问题四都围绕这条主线展开。
